\section{Reproducibilidad del estudio}
\label{cap:apendice}

Este apéndice resume las condiciones necesarias para reproducir el
estudio y relaciona las figuras complementarias que respaldan los
resultados principales. La reproducibilidad se plantea a nivel
metodológico: mismo entorno de ejecución, mismas semillas, mismas
particiones y mismo orden de fases.

\subsection{Condiciones de reproducción}

Los experimentos se ejecutan con Python 3.12 y las bibliotecas
científicas habituales para aprendizaje automático y simulación
numérica: scikit-learn, pandas, numpy, scipy, matplotlib, seaborn,
XGBoost y el SDK de Amazon Braket para la simulación analógica. Todos
los procesos aleatorios se fijan mediante semillas explícitas.

La reproducción completa requiere ejecutar las diez fases en orden. Las
fases 1--7 construyen la referencia clásica: análisis exploratorio,
preprocesado, comprobación posterior al preprocesado, particionado, selección
clásica, modelado y comparación final. Las fases 8--9 ejecutan QFS,
evalúan sus subconjuntos con el mismo protocolo y los integran en la
comparación final. La fase 10 no reentrena modelos ni reejecuta QFS:
lee los artefactos canónicos ya generados y construye el espejo
analítico final de la memoria (régimen, métodos, $k$, $\alpha$, $\beta$,
embebido, métricas, tests, consistencia y síntesis). Cada fase consume
únicamente las salidas validadas de las fases anteriores, por lo que el
conjunto de prueba permanece reservado hasta la evaluación final.

\subsection{Coste computacional}

El bloque clásico se ejecuta en torno a diez minutos en un equipo de
sobremesa convencional. Las fases más costosas son la selección clásica,
por la repetición de doce métodos con tres semillas y permutaciones, y
el modelado, que evalúa 260 configuraciones con cuatro modelos. La fase
cuántica añade el coste de la simulación analógica del barrido de
$\beta$ por dataset y la resolución exacta del QUBO dentro de la
envolvente operativa.

\subsection{Trazabilidad de las cifras}

Cada cifra citada en los capítulos de resultados procede de una salida
tabular generada por la fase correspondiente. La comparación clásica
final, la versión extendida con QFS, la tabla maestra de candidatos, los
contrastes de modelado, la estabilidad de la selección y el control de
calidad del método cuántico se conservan por separado para que cada
afirmación pueda verificarse. El grafo de características que alimenta al
método cuántico, formado por el vector $I_i$ y la matriz $R_{ij}$ de cada
formulación, se genera exclusivamente con datos de entrenamiento y se
mantiene fijo durante la evaluación QFS\@. Las tablas densas de
verificación y caracterización (señal, invarianza del preprocesado,
particiones, perfil de los selectores y cadena de contrastes) se
presentan en el capítulo de resultados, en su punto de uso.


\subsection{Figuras complementarias de verificación}
\label{sec:ap-figuras-auditoria}

Las figuras complementarias se agrupan por función de defensa. Las
primeras documentan calidad de datos y selección clásica; las
siguientes conservan detalle de interpretabilidad y QFS que no debe
sobrecargar el cuerpo de resultados.

\subsubsection{Datos, particiones y selección clásica}
\label{sec:ap-datos-seleccion}

Estas figuras documentan el régimen de cada conjunto antes de QFS\@
(Figura~\ref{fig:ap-regimenes}), la separación frente al azar de la
selección (Figura~\ref{fig:ap-permutacion}) y la pantalla de
\textit{leakage} (Figura~\ref{fig:ap-leakage}).

\begin{figure}[!htbp]
  \centering
  \includegraphics[width=\textwidth]{f10_a_regimenes_dataset.png}
  \caption[Regímenes del dato antes de ejecutar QFS]{Regímenes del dato
    antes de ejecutar QFS\@. La figura reúne efecto supervisado, VIF,
    componentes PCA necesarios para explicar el 80\,\% de la varianza y
    AUC adversarial del particionado para mostrar que QFS no se evalúa
    sobre cinco problemas equivalentes: \textit{Madelon} combina señal
    marginal débil y alta dimensionalidad, \textit{Breast Cancer} y
    \textit{Olive Oil} concentran señal con redundancia estructural, y
    \textit{Customer Churn} opera cerca del techo predictivo.}
  \label{fig:ap-regimenes}
\end{figure}
\FloatBarrier

\begin{figure}[!htbp]
  \centering
  \includegraphics[width=\textwidth]{a1_permutacion_senal_nulo.png}
  \caption[Margen de señal sobre el nulo por variable]{Cuánto supera la
    señal real al azar, variable a variable: cociente entre la puntuación
    real y el percentil 95 de la distribución nula (20 permutaciones),
    para los dos filtros de relevancia sometidos a permutación (F de
    ANOVA e información mutua). La línea vertical en $1$ es el nulo; a su
    derecha la variable supera al azar y la distancia mide la fuerza de
    la señal. Los paneles se ordenan por mediana de señal y eliminan la
    cuadrícula para que se lea la distancia al nulo: \textit{Madelon}
    queda pegado al umbral, mientras que los demás conjuntos muestran un
    margen mucho más holgado.}
  \label{fig:ap-permutacion}
\end{figure}
\FloatBarrier

\begin{figure}[!htbp]
  \centering
  \includegraphics[width=0.86\textwidth]{a2_leakage.png}
  \caption[Pantalla de \textit{leakage}]{Pantalla de \textit{leakage}: AUC univariante y NMI (información mutua normalizada) frente al
    umbral duro de 0{,}99.}
  \label{fig:ap-leakage}
\end{figure}
\FloatBarrier


\subsubsection{Interpretabilidad y deltas de modelado}
\label{sec:ap-interpretabilidad}

La concordancia SHAP media por clase en \textit{Olive Oil}
(Figura~\ref{fig:ap-shap}) y el panorama de deltas de macro-F1 por
selector y dataset (Figura~\ref{fig:ap-deltas}) complementan la lectura
de interpretabilidad del cuerpo.

\begin{figure}[!htbp]
  \centering
  \includegraphics[width=0.86\textwidth]{a4_shap_concordancia.png}
  \caption[Concordancia SHAP por clase en \textit{Olive Oil}]{Importancia SHAP media absoluta por clase en las dos
    formulaciones de \textit{Olive Oil}.}
  \label{fig:ap-shap}
\end{figure}
\FloatBarrier

\begin{figure}[!htbp]
  \centering
  \includegraphics[width=\textwidth]{a5_panorama_deltas.png}
  \caption[Panorama de deltas frente al baseline por selector]{Panorama de
    deltas de macro-F1 frente al baseline (mejor modelo de cada selector,
    en validación), por dataset. La línea vertical en $0$ es el baseline:
    los puntos verdes a su derecha mejoran el baseline y los rojos a su
    izquierda lo empeoran. Cada panel conserva su propia escala porque la
    magnitud de los deltas difiere mucho entre conjuntos (de $\pm0{,}008$
    en \textit{Customer Churn} a $+0{,}31$ en \textit{Madelon}).}
  \label{fig:ap-deltas}
\end{figure}
\FloatBarrier



\subsubsection{Solapes, métricas secundarias y consistencia}
\label{sec:ap-qfs-consistencia}

El solape de QFS-NA con las referencias mRMR y Boruta
(Figura~\ref{fig:ap-solape-qfs}) y la relación entre macro-F1 y AUC-ROC
en las formulaciones binarias (Figura~\ref{fig:ap-auc}) completan las
métricas secundarias.

\begin{figure}[!htbp]
  \centering
  \includegraphics[width=0.86\textwidth]{a8_solape_qfs_clasicos.png}
  \caption[Solape de QFS-NA con mRMR y Boruta frente al azar]{Solape de
    Jaccard entre el subconjunto QFS-NA y las referencias clásicas mRMR
    y Boruta, al presupuesto $k$ más cercano. La banda gris marca el
    solape esperable entre subconjuntos aleatorios del mismo tamaño
    sobre el mismo número de variables: un punto a su derecha indica
    parecido real; un punto dentro de la banda, un solape compatible con
    el azar. Así, el $0{,}62$ de Boruta en \textit{Olive Oil} es trivial
    (Boruta confirma casi todas las variables, de modo que coincidir
    es inevitable), mientras que la coincidencia plena con mRMR en
    \textit{Olive Oil} (3 clases) sí es informativa.}
  \label{fig:ap-solape-qfs}
\end{figure}
\FloatBarrier

\begin{figure}[!htbp]
  \centering
  \includegraphics[width=0.86\textwidth]{a9_macro_f1_auc_binarios.png}
  \caption[Macro-F1 frente a AUC-ROC en binarios]{Macro-F1 frente a AUC-ROC en las formulaciones binarias para
    baseline, QFS-NA y QFS-oráculo. La lectura de apoyo es que el pequeño
    deterioro de \textit{Customer Churn} no es un artefacto exclusivo de
    macro-F1: también se observa en AUC, aunque en un régimen casi
    saturado.}
  \label{fig:ap-auc}
\end{figure}
\FloatBarrier

\begin{figure}[!htbp]
  \centering
  \includegraphics[width=0.86\textwidth]{a7_coste_cuantico.png}
  \caption[Coste de la vía cuántica simulada]{Coste de la vía cuántica simulada: tiempo medio por barrido de
    $\beta$. La figura no busca competir en eficiencia con los filtros
    clásicos, sino cuantificar el precio experimental de explorar QFS-NA
    con simulación analógica en el protocolo de este trabajo.}
  \label{fig:ap-coste-cuantico}
\end{figure}
\FloatBarrier

\begin{figure}[!htbp]
  \centering
  \includegraphics[width=\textwidth]{f10_b9_atomos_mds_error.png}
  \caption[Embebido atómico (MDS) por dataset]{Embebido atómico (MDS)
    verificado por dataset, con el error medio de embebido en cada
    panel. El error de \textit{Customer Churn} es similar al de
    \textit{Breast Cancer} y no es el peor, por lo que la geometría no
    explica por sí sola su deterioro.}
  \label{fig:ap-atomos-mds}
\end{figure}
\FloatBarrier

\begin{figure}[!htbp]
  \centering
  \includegraphics[width=\textwidth]{f10_b2_jaccard_12_metodos.png}
  \caption[Parecido de QFS-NA con los doce selectores frente al azar]{Solape
    de Jaccard entre el subconjunto QFS-NA y cada uno de los doce
    selectores clásicos al presupuesto de referencia (número en cada
    celda). El color no es el Jaccard crudo sino el Jaccard menos el
    solape esperable por azar: verde si QFS se parece a ese selector más
    de lo que lo harían dos subconjuntos aleatorios del mismo tamaño,
    rojo si se parece igual o menos. Así se distingue la afinidad real
    (mRMR en \textit{Olive Oil} de 3 clases, correlación mutua en
    \textit{Customer Churn}) del parecido trivial.}
  \label{fig:ap-jaccard-12}
\end{figure}
\FloatBarrier
